% Unconditional cohomological obstruction; it is not asserted to be the
% obstruction of an ordinary unit-equation pair.
\subsection{A concrete integral localization kernel over \texorpdfstring{\(\mathbb Q\)}{Q}}
\label{sec:ghost-h2-kernel}

The exact localization maps do not force arbitrary higher-weight
integral obstructions to vanish. The following computation gives
a specific nonzero kernel over the same ground field.

\begin{proposition}[The weight-twelve, coefficient-691 kernel]
\label{prop:ghost-h2-691}
Put \(\ell=691\), \(T=\mathbb Z_{691}(12)\), and
\(S=\{691,\infty\}\). Inflation is injective and its image
has the following properties:
\[
 H^2_{\mathrm{cts}}(G_{\mathbb Q,S},T)
       \simeq\mathbb Z/691\mathbb Z
       \lhook\joinrel\longrightarrow H^2_{\mathrm{cts}}(G_{\mathbb Q},T).
\]
Every class in this image restricts to zero in
\(H^2_{\mathrm{cts}}(G_{\mathbb Q_v},T)\) for every place
\(v\). Its image in \(H^2(G_{\mathbb Q},T/691^nT)\) is zero
for every \(n\geq1\).
\end{proposition}
\begin{proof}
The natural Chern-class comparison and the known calculation
of \(K_{22}(\mathbb Z)\) give
\[
 H^2_{\acute et}(\operatorname{Spec}\mathbb Z[1/691],
                  \mathbb Z_{691}(12))
       \simeq K_{22}(\mathbb Z)\{691\}
       \simeq\mathbb Z/691\mathbb Z;
\]
the precise references are \cite[VI.8.2 and VI.10.3]{GhostWeibelKbook}.
The finite-coefficient group-to-scheme maps are isomorphisms by
\cite[II.2.9]{GhostMilneADT}. Their inverse limits identify this
group with \(H^2_{\mathrm{cts}}(G_{\mathbb Q,S},T)\): the
degree-one groups at finite coefficients are finite, so their
derived inverse limit vanishes. This last finiteness follows
also from the bounded-degree, fixed-ramification argument in
Lemma~\ref{lem:ghost-higher-z1}.

We prove injection of inflation directly on continuous cochains.
Let \(c\) be a normalized two-cocycle on \(G_{\mathbb Q,S}\)
whose inflation equals \(\delta f\), with
\(f:G_{\mathbb Q}\to T\) continuous. Let
\(J=\ker(G_{\mathbb Q}\to G_{\mathbb Q,S})\). It acts
trivially on \(T\). Normalization gives \(f(1)=0\); the
coboundary equation with one argument in \(J\) gives
\[
 f(ng)=f(n)+f(g),\qquad
 f(gn)=f(g)+g f(n),\qquad n\in J.
\]
Thus \(f|_J\) is a continuous homomorphism and
\(f(gng^{-1})=g f(n)\). At every \(q\ne691\), wild inertia
and the non-691 parts of tame inertia have zero image. On
691-primary tame inertia, Frobenius conjugation gives
\((q^{12}-q)f(\tau)=0\). Hence all outside-\(S\) inertia
groups are killed. They normally generate \(J\), so \(f|_J=0\).
The same identities show that \(f\) descends continuously to
\(G_{\mathbb Q,S}\), where it is already a primitive of \(c\).
This proves injectivity without interchanging inflation and
an inverse limit of finite-level kernels.

At \(q\ne691\), the restriction of an inflated class factors
through the unramified procyclic quotient. Its continuous
second cohomology with coefficients \(T\) is zero: finite
coefficient cohomology vanishes in degree two, while the
degree-one groups are finite and have zero derived inverse
limit. At \(q=691\), local Tate duality
\cite[I.2.3]{GhostMilneADT} identifies the Pontryagin dual of
the integral second cohomology with
\[
       (\mathbb Q_{691}/\mathbb Z_{691})(-11)^{G_{\mathbb Q_{691}}}.
\]
The cyclotomic image contains every Teichm\"uller unit of order
dividing 690. Since \(690\nmid11\), one of these acts on the
displayed module by a scalar whose difference from one is a
691-adic unit. The invariants are zero. At the real place,
averaging over the group of order two kills positive cohomology
with these odd-primary coefficients. This proves every asserted
integral local vanishing.

For completeness we check exactly why all absolute finite
coefficient reductions vanish too. Put
\(T_n=T/691^nT\); its Tate dual is
\(T_n^D=(\mathbb Z/691^n\mathbb Z)(-11)\). Its splitting
field is contained in \(\mathbb Q(\mu_{691^n})\), whose Galois
group is cyclic. The localization map on its first cohomology
is injective. Indeed, for a locally trivial class, restriction
to its splitting field is a homomorphism to a finite constant
module, locally zero everywhere; Chebotarev makes this
homomorphism zero. The class then inflates from the finite
cyclic splitting group. A prime with Frobenius generating
that group gives the full group as a decomposition group;
local triviality and the injective local inflation map make
the inflated class zero there and hence zero globally.
Poitou--Tate duality, in the all-places form of
\cite[I.4.10]{GhostMilneADT}, identifies the kernel of
localization on \(H^2(G_{\mathbb Q},T_n)\) with the dual
of this zero first-cohomology kernel. The reduction of our
class has zero localizations, so it is zero.
\end{proof}

\begin{corollary}[An actual continuous extension with every finite reduction split]
\label{cor:ghost-integral-extension}
There is an exact sequence of profinite groups
\[
  1\longrightarrow\mathbb Z_{691}(12)
    \longrightarrow\mathcal E\longrightarrow G_{\mathbb Q}
    \longrightarrow1
\]
with no continuous group-theoretic section, such that its
restriction to every \(G_{\mathbb Q_v}\) has a continuous
section and \(\mathcal E/691^nT\to G_{\mathbb Q}\) has a
continuous section for every \(n\geq1\).
\end{corollary}
\begin{proof}
Choose a continuous normalized two-cocycle \(c\) representing
a nonzero inflated class of Proposition~\ref{prop:ghost-h2-691}.
On the compact space \(T\times G_{\mathbb Q}\) define
\[
 (u,g)(v,h)=(u+g v+c(g,h),gh).
\]
The cocycle identity is precisely associativity. This is a
compact totally disconnected topological group, hence a
profinite group, and its projection gives the displayed
extension. A continuous section has the form \(g\mapsto(b(g),g)\)
and is a homomorphism exactly when \(\delta b=-c\).
The nonzero class therefore excludes a global section. The
proved local and finite-coefficient vanishings provide exactly
such continuous primitives for each restriction and each
reduction, and hence give all the stated sections.
\end{proof}

This extension is defined by the displayed cochain formula.
No map identifying it with a quotient of
\(\pi_1(\mathbb P^1_{\mathbb Q}\setminus\{0,1,\infty\})\)
has been constructed here. The computation establishes the
failure of the proposed general higher-weight localization
injection; it does not evaluate a particular ordinary pair's
nilpotent obstruction. The fixed-pair compact-tower theorem
does not require that injection and retains its conclusion.

\subsection{An uncountable localization kernel already in weight two}
\label{sec:ghost-uncountable-kernel}

\begin{proposition}[The exact derived-limit kernel and a regular-prime example]
\label{prop:ghost-uncountable-kernel}
For every odd prime \(\ell\) and every \(r\geq2\), the
kernel of the actual integral localization map is exactly
\[
 \ker\!\left(H^2_{\rm cts}(G_{\mathbb Q},\mathbb Z_\ell(r))
       \longrightarrow\prod_v H^2_{\rm cts}
                   (G_{\mathbb Q_v},\mathbb Z_\ell(r))\right)
 \simeq\varprojlim_n^{\,1}H^1
                (G_{\mathbb Q},\mathbb Z/\ell^n(r)).
\]
At \(\ell=3,r=2\), this kernel contains a subgroup isomorphic
to
\[
 \frac{\prod_{j\geq1}\mathbb Z_3}
      {c_0(\mathbb Z_3)},\qquad
 c_0(\mathbb Z_3)=
 \{(z_j):z_j\longrightarrow0\text{ in }\mathbb Z_3\}.
\]
In particular it is uncountable. The unramified integral group
\(H^2(G_{\mathbb Q,\{3,\infty\}},\mathbb Z_3(2))\)
is nevertheless zero.
\end{proposition}
\begin{proof}
For the coefficient tower \(T_n=\mathbb Z/\ell^n(r)\), the
continuous cochain complexes have surjective transition maps
in every degree: a continuous finite-valued cochain lifts by
choosing a lift of each of its finitely many values. The
short exact sequence of complexes defined by one minus the
shift on their product yields
\[
 0\longrightarrow\varprojlim{}^1 H^1(G,T_n)
 \longrightarrow H^2_{\rm cts}(G,\mathbb Z_\ell(r))
 \longrightarrow\varprojlim H^2(G,T_n)\longrightarrow0.
\]
This follows directly by taking cohomology: the kernel and
cokernel of one minus the shift on \(\prod_n H^i(G,T_n)\)
are the definitions of \(\varprojlim\) and
\(\varprojlim{}^1\). For a tower of finite groups the latter
cokernel is zero, since every finite collection of equations
for the shift has a solution, and compactness of the product
of the finite groups gives a simultaneous solution.

The finite dual \(T_n^D=\mathbb Z/\ell^n(1-r)\) has cyclic
splitting group, because \(\ell\) is odd. The Chebotarev
and Poitou--Tate argument in
Proposition~\ref{prop:ghost-h2-691} therefore proves that
every finite-coefficient global \(H^2\) localization is
injective. At each place, finite-coefficient \(H^1\) is
finite \cite[I.2.1]{GhostMilneADT}, so the local derived
limit is zero. An integral global class is locally zero
exactly when all its absolute finite reductions are zero.
The displayed exact sequence proves the first assertion.

Now put \(\ell=3,r=2\), \(S=\{3,\infty\}\), and
\(X=\operatorname{Spec}\mathbb Z[1/3]\). The comparison
already used above gives
\[
 H^2(G_{\mathbb Q,S},\mathbb Z_3(2))
      \simeq K_2(\mathbb Z)\{3\}=0
 \quad\cite[VI.8.2, VI.10.1.1]{GhostWeibelKbook}.
\]
Positive-degree finite cohomology above degree two vanishes
for this odd coefficient prime
\cite[I.4.10(c)]{GhostMilneADT}. Passing to integral
cohomology preserves degree-three vanishing, because the
finite degree-two groups are finite and have zero derived
limit. The coefficient sequence therefore gives
\(H^2(G_{\mathbb Q,S},T_n)=0\) for every \(n\).

The exact finite-coefficient residue sequence is
\[
 0\longrightarrow A_n
   \longrightarrow B_n
   \xrightarrow{\partial_n}D_n\longrightarrow0,
 \quad
 \begin{cases}
 A_n=H^1(G_{\mathbb Q,S},T_n),\\
 B_n=H^1(G_{\mathbb Q},T_n),\\
 D_n=\displaystyle\bigoplus_{q\ne3}
                   H^0(\mathbb F_q,\mathbb Z/3^n(1)).
 \end{cases}
\]
Here \(\partial_n\) is the restriction to inertia, followed
by the tame-character identification. To justify precisely
the maps and exactness, remove finitely many closed points
from \(X\), use the cohomology sequence with supports, and
then take the filtered union over those finite sets. At a
point \(q\ne3\), purity gives zero support cohomology in
degree one and
\(H^2_q(X,T_n)=H^0(\mathbb F_q,T_n(-1))\).
The needed purity calculation is
\cite[II.1.7(b)]{GhostMilneADT}: its formula
\(R^2i^!\mu_{3^n}=\mathbb Z/3^n\), with the other terms
zero, tensors with the locally free sheaf \(\mu_{3^n}\)
to give exactly this twist. The group-to-scheme comparison
is II.2.9 of the same source. The next global term is the
zero group \(H^2(X,T_n)\), proving surjectivity of the
displayed residue map. All these maps commute with standard
coefficient reduction; the residue terms have finite support
because each finite-valued Galois cocycle is ramified at
only finitely many primes.

Each \(A_n\) is finite by fixed-ramification finiteness.
The inverse-limit exact sequence and the vanishing of its
derived limit give
\(\varprojlim{}^1 B_n\simeq\varprojlim{}^1D_n\).
For an explicit direct summand, choose distinct primes
\(q_j\equiv1+3^j\pmod{3^{j+1}}\), \(j\geq1\).
Chebotarev in the cyclotomic field supplies them
\cite[Theorem~8.31]{GhostMilneANT}, and
\(v_3(q_j-1)=j\). Put \(d_j=q_j-1\).
Choosing compatible Tate bases identifies the selected
summand of \(D_n\) with
\[
 D'_n=\bigoplus_j
   \ker(d_j:\mathbb Z/3^n\longrightarrow\mathbb Z/3^n).
\]
Projection onto the selected primes splits the inclusion
as an inverse system, so \(\varprojlim{}^1D'_n\) is a
direct summand of \(\varprojlim{}^1D_n\).

Consider the two-term complexes
\[
 P_n^\bullet=\left[
 \bigoplus_j\mathbb Z/3^n\xrightarrow{(d_j)}
 \bigoplus_j\mathbb Z/3^n\right]
\]
in degrees zero and one. Their transition maps are
degreewise surjective, their degree-zero cohomology is
\(D'_n\), and their inverse limit is
\([R\xrightarrow{d}R]\), where \(R=c_0(\mathbb Z_3)\).
The same product-complex argument gives
\[
 0\longrightarrow\varprojlim{}^1D'_n
 \longrightarrow R/dR
 \longrightarrow\varprojlim_n H^1(P_n^\bullet)
 \longrightarrow0.
\]
The last map reduces each coordinate modulo \(d_j\).
Consequently its kernel consists exactly of the classes
of sequences \(z\in R\) with \(z_j\in d_j\mathbb Z_3\)
for every \(j\). Since \(v_3(d_j)=j\to\infty\),
multiplication by \(d\) maps the full product
\(\prod_j\mathbb Z_3\) into \(R\), injectively. Every
element of this kernel is \(dy\) for a unique such
sequence \(y\), and it belongs to \(dR\) exactly when
\(y\in R\). Thus
\[
 \varprojlim{}^1D'_n\simeq
       (\prod_j\mathbb Z_3)/c_0(\mathbb Z_3).
\]
Two binary sequences have the same class here exactly
when they differ at finitely many coordinates. There are
continuum many resulting classes, which proves
uncountability and the claimed explicit injection.
\end{proof}

The proof identifies the exact coefficient maps, residue
maps, and derived limits. It does not assert that any
displayed nonzero class equals the obstruction of an
ordinary pair. In particular ordinary unit-equation
relations may force special obstruction representatives
to vanish even in this large ambient kernel; that is a
separate calculation at each nilpotency class.
